\chapter{Complexity — Parse Plain}
%%% generated by the blueprint pipeline from TCSlib.Complexity.SpaceComplexity.Machines.ParsePlain
\section{Overview}

Declarations of the Lean module \texttt{TCSlib.Complexity.SpaceComplexity.Machines.ParsePlain}.

\section{Declarations}

\begin{definition}[Valid plain encoding]
\lean{Complexity.LogProg.ValidPlain}
\label{Complexity.LogProg.ValidPlain}
\leanok
\uses{Complexity.LogProg.Canon, Turing.pairEncode}
A Boolean list $y$ is a valid plain encoding if there are a natural number $n$ and a
list $w$ such that $y$ is the pair encoding of the list of $n$ copies of true with $w$,
and $w$ is canonical.
\end{definition}

\begin{definition}[Valid pair encoding]
\lean{Complexity.LogProg.ValidPair}
\label{Complexity.LogProg.ValidPair}
\leanok
\uses{Complexity.LogProg.Canon, Turing.pairEncode}
A Boolean list $y$ is a valid pair encoding if there are a natural number $n$ and
Boolean lists $u,w$ such that $y$ is the pair encoding of the list of $n$ copies of true
with the pair encoding of $u$ and $w$, and both $u$ and $w$ are canonical.
\end{definition}

\begin{lemma}[Doubling a unary word]
\lean{Complexity.LogProg.dbl_replicate}
\label{Complexity.LogProg.dbl_replicate}
\leanok
\uses{Turing.dbl}
\difficulty{4}
For every natural number $n$, the bitwise doubling of the word consisting of $n$ copies
of $1$ is the word consisting of $2n$ copies of $1$.
\end{lemma}

\begin{definition}[Updated machine configuration]
\lean{Complexity.LogProg.xCfg}
\label{Complexity.LogProg.xCfg}
\leanok
\uses{Turing.Cfg}
Given a configuration of an $m$-work-tape Turing machine with Boolean input $x$, a state
$s$, an input-head position $i$, a work-tape index $r$, and an integer position $p$,
form a configuration whose state is $s$, input-head position is $i$, work tapes and
output are unchanged, and whose work-tape head positions agree with the original except
that tape $r$ has head position $p$.
\end{definition}

\begin{lemma}[State of Updated Configuration]
\lean{Complexity.LogProg.xCfg_state}
\label{Complexity.LogProg.xCfg_state}
\leanok
\uses{Complexity.LogProg.xCfg, Turing.Cfg}
\difficulty{1}
Let $m$ be a natural number, let $x$ be a finite bit string, and let $\Lambda$ be a
type. Given any configuration $c$ for input $x$ with $m$ work tapes, any state
$s\in\Lambda$, any valid input-head position $ip$, any work-tape index $r<m$, and any
integer position $p$, the state of the updated configuration is $s$.
\end{lemma}

\begin{definition}[Work-tape movement action]
\lean{Complexity.LogProg.xAct}
\label{Complexity.LogProg.xAct}
\leanok
\uses{Turing.Action}
For an $m$-work-tape machine, a work-tape index $r$, input-head movement $a$, work-tape
movement $b$, and state $s$, define an action that attempts movement $a$ on the input
head, makes no writes on any work tape, moves tape $r$ by $b$ and leaves the other
work-tape heads stationary, emits no output symbol, and enters state $s$.
\end{definition}

\begin{definition}[Rejecting action]
\lean{Complexity.LogProg.rejAct}
\label{Complexity.LogProg.rejAct}
\leanok
\uses{Turing.Action}
Define an action that leaves the input head and all work-tape heads stationary, makes no
work-tape writes, emits false on the output tape, and halts.
\end{definition}

\begin{lemma}[Moving the input and a register head]
\lean{Complexity.LogProg.apply_xAct}
\label{Complexity.LogProg.apply_xAct}
\leanok
\uses{Complexity.LogProg.xAct, Complexity.LogProg.xCfg, Turing.Action.apply, Turing.Cfg, Turing.moveInputPos}
\difficulty{5}
Let $c$ be a configuration with $m$ Boolean work tapes, state set $\Lambda$, and fixed
input $x$. Given states $s,s'\in\Lambda$, an input-head position
$ip\in\mathrm{Fin}(|x|+2)$, a register $r\in\mathrm{Fin}(m)$, an integer $p$, and input
and register movements $mv_I,mv_R\in\{-1,0,1\}$, applying the action that moves those
heads by the respective movements and enters $s'$ to the configuration obtained from $c$
by setting its state to $s$, input-head position to $ip$, and register-head position to
$p$ gives the configuration with state $s'$, input-head position obtained by the clamped
move from $ip$ by $mv_I$, and register-head position $p+mv_R$. All other configuration
data are unchanged.
\end{lemma}

\begin{lemma}[Work-Tape Symbol After Configuration Update]
\lean{Complexity.LogProg.xCfg_read}
\label{Complexity.LogProg.xCfg_read}
\leanok
\uses{Complexity.LogProg.xCfg, Turing.Cfg, Turing.Cfg.workTapeSymbols}
\difficulty{1}
Let $m$ be a natural number, let $x$ be a finite bit string, and let $\Lambda$ be a
type. Given any configuration $c$ for input $x$ with $m$ work tapes, any state
$s\in\Lambda$, any valid input-head position $ip$, any work-tape index $r<m$, and any
integer position $p$, the symbol read by work tape $r$ in the updated configuration
equals the content of tape $r$ at position $p$ in $c$.
\end{lemma}

\begin{definition}[Input symbol at a shifted position]
\lean{Complexity.LogProg.inSym}
\label{Complexity.LogProg.inSym}
\leanok
For a Boolean input list $x$ and a natural number $q$, return no symbol when $q=0$;
otherwise return the Boolean at list index $q-1$, if that index lies within $x$, and no
symbol if it does not.
\end{definition}

\begin{lemma}[Input Symbol After Configuration Update]
\lean{Complexity.LogProg.xCfg_inSym}
\label{Complexity.LogProg.xCfg_inSym}
\leanok
\uses{Complexity.LogProg.inSym, Complexity.LogProg.inputSymbol_eq, Complexity.LogProg.xCfg, Turing.Cfg, Turing.Cfg.inputSymbol}
\difficulty{3}
Let $m$ be a natural number, let $x$ be a finite bit string, and let $\Lambda$ be a
type. Given any configuration $c$ for input $x$ with $m$ work tapes, any state
$s\in\Lambda$, any input-head position $ip$ among the $|x|+2$ positions, any work-tape
index $r<m$, and any integer position $p$, the updated configuration reads the same
input symbol as the symbol at position $ip$ in $x$, with the two boundary positions
interpreted as blank.
\end{lemma}

\begin{lemma}[One ordinary transition from an input configuration]
\lean{Complexity.LogProg.rrun_one_x}
\label{Complexity.LogProg.rrun_one_x}
\leanok
\uses{Complexity.LogProg.inSym, Complexity.LogProg.rrun, Complexity.LogProg.rrun_one, Complexity.LogProg.rstep_noncall, Complexity.LogProg.xCfg, Complexity.LogProg.xCfg_inSym, Turing.Action.apply, Turing.Cfg, Turing.Cfg.workTapeSymbols, Turing.MultiTapeTM.step}
\difficulty{3}
Let $P$ be a register program with oracle on Boolean input $x$, and let $c$ be a
configuration. Set its state to $s$, its input-head position to $ip$, and register $r$'s
head position to $p$, keeping its tape contents, other head positions, and output from
$c$. If $P$ has no oracle call at $s$, then one step of its run from this configuration
equals applying the machine transition prescribed by $P$ for state $s$, the input symbol
at position $ip$ (with no symbol at position zero), and the scanned work-tape symbols.
\end{lemma}

\begin{lemma}[Rightward input-head movement]
\lean{Complexity.LogProg.moveInputPos_pos_val}
\label{Complexity.LogProg.moveInputPos_pos_val}
\leanok
\uses{Complexity.LogProg.moveInputPos_val, Turing.moveInputPos}
\difficulty{3}
For every Boolean input word $x$ and every input-head position $ip$ among the $|x|+2$
positions whose value is at most $|x|$, moving the head one step to the right gives
position value $ip+1$.
\end{lemma}

\begin{lemma}[Leftward input-head movement]
\lean{Complexity.LogProg.moveInputPos_neg_val'}
\label{Complexity.LogProg.moveInputPos_neg_val'}
\leanok
\uses{Complexity.LogProg.moveInputPos_val, Turing.moveInputPos}
\difficulty{1}
For every Boolean input word $x$ and every input-head position $ip$ among the $|x|+2$
positions, moving the head one step to the left gives position value $ip-1$.
\end{lemma}

\begin{lemma}[Simulation of a scan across input symbols]
\lean{Complexity.LogProg.scanR}
\label{Complexity.LogProg.scanR}
\leanok
\uses{Complexity.LogProg.apply_xAct, Complexity.LogProg.inSym, Complexity.LogProg.moveInputPos_pos_val, Complexity.LogProg.rrun, Complexity.LogProg.rrun_one, Complexity.LogProg.rrun_one_x, Complexity.LogProg.rrun_succ, Complexity.LogProg.xAct, Complexity.LogProg.xCfg, Turing.Cfg}
\difficulty{4}
Let $P$ be a register program with $m$ registers, an oracle, an input word $x$, and a
configuration $c$; fix a register $r$, an integer position $p$, a state sequence
$\sigma_j$, and natural numbers $q_0,k$. Suppose $q_0+k\le |x|+1$. Suppose that for
every $j<k$ and every tuple of scanned register symbols, the ordinary transition from
state $\sigma_j$, reading the input symbol at position $q_0+j$, is the action that moves
the input head right, leaves every register tape unchanged and stationary, emits no
output, and enters state $\sigma_{j+1}$. Suppose also that there is no oracle call at
any state $\sigma_j$ for $j<k$. Then for every $j\le k$, after $j$ steps from the
configuration with state $\sigma_0$, input-head position $q_0$, and register-$r$ head
position $p$, the run has state $\sigma_j$, input-head position $q_0+j$, and
register-$r$ head position $p$, with the other configuration data unchanged.
\end{lemma}

\begin{lemma}[Rewinding the input head]
\lean{Complexity.LogProg.rewind_x}
\label{Complexity.LogProg.rewind_x}
\leanok
\uses{Complexity.LogProg.apply_xAct, Complexity.LogProg.inSym, Complexity.LogProg.moveInputPos_neg_val', Complexity.LogProg.moveInputPos_pos_val, Complexity.LogProg.rrun, Complexity.LogProg.rrun_add, Complexity.LogProg.rrun_one_x, Complexity.LogProg.rrun_succ_left, Complexity.LogProg.xAct, Complexity.LogProg.xCfg, Turing.Cfg, Turing.moveInputPos}
\difficulty{6}
Let $P$ be a register program with $m$ registers and $d$ deciders, let $x$ be a Boolean
input, and let $\mathcal O$ be an oracle assigning a Boolean answer to each decider and
Boolean word. Suppose the ordinary transitions at states $rw_1$ and $rw_2$ always act as
follows: at $rw_1$ the input head moves left and the machine enters $rw_2$; at $rw_2$,
it moves left and stays in $rw_2$ when the input symbol is present, and moves right and
enters $nx$ when the input symbol is absent. Suppose neither $rw_1$ nor $rw_2$ is a call
state. From a configuration whose state is $rw_1$, whose input-head position is $ip$,
and whose register $r$ head is at integer position $p$, there is a time $T$ when the
register head is still at $p$, the state is $nx$, and the input-head position is $1$. At
every earlier time $t<T$, the configuration has some state and input-head position, the
register $r$ head is still at $p$, and that state is not a call state.
\end{lemma}

\begin{definition}[Initial true-prefix length]
\lean{Complexity.LogProg.tRun}
\label{Complexity.LogProg.tRun}
\leanok
For a Boolean list $x$, define its initial true-prefix length to be the number of
consecutive true entries at the beginning of $x$.
\end{definition}

\begin{lemma}[Bound on the initial run of true bits]
\lean{Complexity.LogProg.tRun_le}
\label{Complexity.LogProg.tRun_le}
\leanok
\uses{Complexity.LogProg.tRun}
\difficulty{1}
For every finite Boolean list $x$, the length of its initial run of true entries is at
most the length of $x$.
\end{lemma}

\begin{lemma}[Bits in the Initial Run of Ones]
\lean{Complexity.LogProg.getElem?_lt_tRun}
\leanok
\uses{Complexity.LogProg.tRun, Complexity.LogProg.takeWhile_true_getElem?}
\difficulty{1}
For a Boolean word $x$, let $t$ be the length of its initial run of ones. Every index
$j<t$ has entry $x_j=\mathrm{true}$.
\end{lemma}

\begin{lemma}[The Initial Run Ends at a Non-One Entry]
\lean{Complexity.LogProg.getElem?_tRun}
\leanok
\uses{Complexity.LogProg.tRun, Complexity.LogProg.takeWhile_true_end}
\difficulty{1}
For a Boolean word $x$, let $t$ be the length of its initial run of ones. The possibly
undefined entry of $x$ at index $t$ is not true.
\end{lemma}

\begin{lemma}[Characterization of valid unary-prefixed words]
\lean{Complexity.LogProg.validPlain_iff}
\label{Complexity.LogProg.validPlain_iff}
\leanok
\uses{Complexity.LogProg.Canon, Complexity.LogProg.ValidPlain, Complexity.LogProg.dbl_replicate, Complexity.LogProg.getElem?_lt_tRun, Complexity.LogProg.tRun, Complexity.LogProg.tRun_le, Turing.pairEncode, Turing.pairEncode_eq_dbl}
\difficulty{6}
For every Boolean word $y$, it is valid plain if and only if its initial run of ones has
even length, the next two bits are $01$, and the suffix after those two bits is
canonical (empty or ending in $1$).
\end{lemma}

\begin{definition}[Validation action]
\lean{Complexity.LogProg.valUAct}
\label{Complexity.LogProg.valUAct}
\leanok
\uses{Complexity.LogProg.rejAct, Complexity.LogProg.xAct, Turing.Action}
For an $m$-work-tape machine, a work-tape index $r$, states $v_0,v_1,v_S$, a parity bit
$e$, and an optional Boolean $a$, define an action as follows. If $a$ is true, move the
input head one position to the right, leave all work-tape heads stationary without
writing, emit no output, and enter $v_0$ when $e$ is true or $v_1$ when $e$ is false. If
$a$ is false and $e$ is true, leave both the input and work-tape heads stationary, make
no writes, emit false, and halt; if $a$ is false and $e$ is false, move the input head
one position to the right, leave all work-tape heads stationary without writing, emit no
output, and enter $v_S$. If $a$ is absent, leave both the input and work-tape heads
stationary, make no writes, emit false, and halt.
\end{definition}

\begin{definition}[State validation action]
\lean{Complexity.LogProg.valSAct}
\label{Complexity.LogProg.valSAct}
\leanok
\uses{Complexity.LogProg.rejAct, Complexity.LogProg.xAct, Turing.Action}
For an input word $x$ and a register index $r$ among $m$ registers, the state-validation
action with parameter $vW0$ returns the action $xAct(r,1,0,vW0)$ when the observed
symbol is $1$, and returns the rejection action for either no symbol or the symbol $0$.
\end{definition}

\begin{definition}[Work-tape validation action]
\lean{Complexity.LogProg.valWAct}
\label{Complexity.LogProg.valWAct}
\leanok
\uses{Complexity.LogProg.rejAct, Complexity.LogProg.xAct, Turing.Action}
For an input word $x$, a register index $r$ among $m$ registers, states $vWF,vWT,rw_1$,
and a remembered bit that may be absent, the action on an observed symbol $a$ is as
follows. If $a$ is a bit $b$, it is $xAct(r,1,0,vWT)$ when $b=1$ and $xAct(r,1,0,vWF)$
when $b=0$. If $a$ is absent, it is the rejection action when the remembered bit is $0$,
and otherwise it is $xAct(r,0,0,rw_1)$.
\end{definition}

\begin{lemma}[Input symbol at a successor position]
\lean{Complexity.LogProg.inSym_succ}
\label{Complexity.LogProg.inSym_succ}
\leanok
\uses{Complexity.LogProg.inSym}
\difficulty{1}
For every Boolean word $x$ and natural number $j$, the input symbol at position $j+1$ is
the $j$th symbol of $x$, with the result undefined when that index is outside the word.
\end{lemma}

\begin{lemma}[One-step rejection behavior]
\lean{Complexity.LogProg.rrun_one_rej}
\label{Complexity.LogProg.rrun_one_rej}
\leanok
\uses{Complexity.LogProg.RProg, Complexity.LogProg.inSym, Complexity.LogProg.rejAct, Complexity.LogProg.rrun, Complexity.LogProg.rrun_one_x, Complexity.LogProg.xCfg, Turing.Action.apply, Turing.Cfg}
\difficulty{3}
Let $P$ be a register program with oracle on Boolean input $x$. Choose a configuration
$c$, a state $s$, an input-head position $ip$, a register $r$, and an integer $p$. Form
a configuration with state $s$, input-head position $ip$, and register $r$'s head at
$p$, leaving the other components as in $c$. Suppose there is no oracle call at $s$, and
for every tuple of scanned work-tape symbols the transition at $s$, on the input symbol
read at position $ip$ (with no symbol at position zero), is the rejection action, which
emits false and halts without moving any work-tape head. Then after one step the state
is halted, the output is $c$'s output followed by false, and the work-tape head
positions are those of $c$ with register $r$'s position replaced by $p$.
\end{lemma}

\begin{definition}[Valid fragment execution]
\lean{Complexity.LogProg.FragOK}
\label{Complexity.LogProg.FragOK}
\leanok
\uses{Complexity.LogProg.RProg, Complexity.LogProg.rrun, Complexity.LogProg.xCfg, Turing.Cfg}
Let $P$ be a register program with $m$ registers and $d$ deciders, let $oracle$ assign a
Boolean answer to each decider and Boolean word, let $c$ be a machine configuration on
input $x$, and choose an initial state $s_0$, a register $r$, a tape position
$p\in\mathbb Z$, a proposition $good$, and a next state $next$. The fragment is valid if
both of the following implications hold: if $good$, there is a time $T$ at which the run
from the configuration formed by setting the state to $s_0$, the input-head position to
$1$, and register $r$'s work-head position to $p$ reaches the corresponding
configuration with state $next$; at every earlier time $t<T$, the run is in a
configuration of that same form for some state and input-head position, and that state
is not a call state. If $good$ is false, there is a time $T$ at which that run has
halted, its output is the original output followed by $0$, and its work-head positions
are the original positions except that register $r$'s position is $p$; at every earlier
time, the same form of configuration and absence of a call state hold.
\end{definition}

\begin{lemma}[Validation of a plain input]
\lean{Complexity.LogProg.valPlain_run}
\label{Complexity.LogProg.valPlain_run}
\leanok
\uses{Complexity.LogProg.FragOK, Complexity.LogProg.ValidPlain, Complexity.LogProg.apply_xAct, Complexity.LogProg.getElem?_lt_tRun, Complexity.LogProg.getElem?_tRun, Complexity.LogProg.inSym, Complexity.LogProg.inSym_succ, Complexity.LogProg.moveInputPos_pos_val, Complexity.LogProg.rejAct, Complexity.LogProg.rewind_x, Complexity.LogProg.rrun, Complexity.LogProg.rrun_add, Complexity.LogProg.rrun_one_rej, Complexity.LogProg.rrun_one_x, Complexity.LogProg.scanR, Complexity.LogProg.tRun, Complexity.LogProg.tRun_le, Complexity.LogProg.valSAct, Complexity.LogProg.valUAct, Complexity.LogProg.valWAct, Complexity.LogProg.validPlain_iff, Complexity.LogProg.xAct, Complexity.LogProg.xCfg, Turing.Cfg}
\difficulty{7}
Let $P$ be a register program with oracle $\mathcal O$, $r$ a register, $p\in\mathbb Z$,
and $c$ a configuration on input $x$. Suppose the transition identities
$hU0,hU1,hS,hW0,hWF,hWT,h_1,h_2$ hold at the specified states and none of those states
is an oracle-call state, as asserted by $cU0,cU1,cS,cW0,cWF,cWT,c_1,c_2$. Then the run
from state $vU0$, input-head position $1$, and register-head position $p$ satisfies the
following: if $x$ is a valid plain encoding, it reaches state $next$ at input-head
position $1$; if $x$ is not a valid plain encoding, it halts with a rejecting output. In
either case, the run preserves the indicated register-head positions and avoids
oracle-call states before its endpoint.
\end{lemma}
